\documentclass[../main]{subfiles}
\begin{document}
\section{Pérdidas en Motores.}
\begin{center}
  \img{images/09/placamotor.jpg}{Placa características de motor}
\end{center}
\info{¿Qué son las pérdidas?}{Las pérdidas de los motores eléctricos
  se producen en las láminas del estator, las láminas del rotor y los
  devanados, pero también en los imanes permanentes. Cuanto menores
sean las pérdidas del motor eléctrico, mayor será su eficiencia, por supuesto.}
\section{Cálculo de pérdidas}
\subsection{Pérdidas en el estator}
\subsubsection{Potencia eléctrica en la entrada del estator.}
\begin{equation}
  P= m \cdot V \cdot I \cos(\varphi)
\end{equation}

siendo:

$P$: Potencia de entrada en el estator (potencia electrica absorvida de la red).

$m$: número de fases del estator.

$V$: Tensión por fase del estator.

$I$: Corriente del estator por fase.

$\cos(\varphi)$: f.d.p del motor (f.d.p.del estator);

\subsubsection{Pérdidas en el cobre del estator ($P_{Cu}$)}
\begin{equation}
  P_{Cu}= m \cdot R \cdot I^2
\end{equation}
\subsubsection{Pérdidas en el estator.}
\begin{equation}
  P_p= P_{Cu} + P_{Fe}
\end{equation}
\subsection{Pérdidas en el hierro del estator}
\begin{equation}
  P_{Fe}=m \cdot E \cdot I_{Fe}
\end{equation}
$I_{Fe}$: Corriente de pérdidas en el hierro.
\subsection{Potencia que atraviesa el entrehierro}
\begin{equation}
  P_a = P- P_p = P-P_{Cu} - P_{Fe}
\end{equation}
siendo

$P$: Potencia de entrada al estator.

$P_p$: Pérdidas en el estator.

\subsection{Pérdidas en el rotor}
\subsubsection{Pérdidas en el cobre del rotor}
\begin{equation}
  P_{Cu2}= m \cdot R I^2
\end{equation}
\subsection{Potencia mecánica interna del motor}
\begin{equation}
  P_m=P_a-P_{Cu2}
\end{equation}
\subsection{Expresión de la potencia mecánica interna del motor}
\begin{equation}
  P_m=m \cdot R_2 (1-s) I_2^2
\end{equation}
\subsection{Potencia mecánica útil de motor}
\begin{equation}
  P_u=P_i-P_m
\end{equation}
\subsection{Rendimiento del motor}
\begin{equation}
  \eta = \dfrac{P_u}{P}=\dfrac{P_u}{P_u+P_m+P_{cu_2}+P_{Fe}+P_{Cu_1}}
\end{equation}
\actividad{Problemas}
\begin{enumerate}
  \item La potencia absorbida por un motor asíncrono trifásico de 4
    polos, 50 Hz, es de 476kW cuando gira a 1435 R.P.M. Las pérdidas
    totales en el estator son de 265W y las de rozamiento y
    ventilación son de 300W.
    Calcular
    \begin{enumerate}
      \item el deslizamiento
      \item las pérdidas en el rotor
      \item potencia útil en el árbol del motor
      \item rendimiento.
    \end{enumerate}
  \item Un motor de inducción trifásico de 8 polos, 1CV, 380V gira a
    72R.P.M a plena carga. Si el rendimiento y el f.d.p. a esta carga
    esd e 83\% y 0,75 respectivamente. Calcular
    \begin{enumerate}
      \item velocidad de sincronismo del campo giratorio.
      \item deslizamiento a plena carga;
      \item corriente de línea
        \par en el árbol de la máquina.
    \end{enumerate}
  \item Un motor asincrono trifásico de 4 polos, conectado en
    estrella se alimenta por una red de 380V, 50Hz. La impedancia del
    estátor es igual a $z=0,1 +j0,5\Omega$ y la del rotor en reposo
    reducida al estator vale $z=0,1+j0,3\Omega/fase$ Calcular
    \begin{enumerate}
      \item intensidad absorbida en el arranque.
      \item corrente a plena carga, si el deslizamiento es de 4\%.
      \item potencia y par nominal si se desprecian las pérdidas mecánicas.
      \item rendimiento en el caso anterior si las pérdidas en el
        hierro son iguales a 1200W.
    \end{enumerate}
\end{enumerate}
\actividad{Problemas}
\begin{enumerate}
    \pagestyle{empty}
    \setcounter{enumi}{3}
  \item Un motor de inducción trifásico de 6 polos, 60 Hz, absorbe
    una potencia de 200 kW cuando gira a 1180 R.P.M. Las pérdidas en
    el estator son de 180 W y las pérdidas por fricción son de 250 W. Calcular:
    \begin{enumerate}
      \item el deslizamiento,
      \item las pérdidas en el rotor,
      \item la potencia útil en el eje del motor,
      \item el rendimiento.
    \end{enumerate}
  \item Un motor asincrónico de 10 polos, 1.5 CV, 440 V, opera a 58
    R.P.M. a plena carga. El rendimiento es del 85\% y el factor de
    potencia es de 0.78. Calcular:
    \begin{enumerate}
      \item la velocidad sincrónica,
      \item el deslizamiento a plena carga,
      \item la corriente de línea,
      \item la potencia entregada en el eje del motor.
    \end{enumerate}
  \item Un motor de inducción trifásico de 4 polos, 50 Hz, conectado
    en triángulo a una red de 400 V tiene una impedancia del estator
    de $z = 0.12 + j0.6 \Omega$. La impedancia del rotor en reposo es
    $z = 0.15 + j0.35 \Omega$. Calcular:
    \begin{enumerate}
      \item la corriente absorbida en el arranque,
      \item la corriente a plena carga con un deslizamiento del 3\%,
      \item el par nominal suponiendo despreciables las pérdidas mecánicas,
      \item el rendimiento si las pérdidas en el hierro son de 1500 W.
    \end{enumerate}
  \item Un motor asíncrono de 8 polos, 1 CV, 380 V, tiene un
    deslizamiento de 2\% cuando gira a plena carga. Las pérdidas
    totales en el estator y rotor son de 500 W. Calcular:
    \begin{enumerate}
      \item la velocidad sincrónica,
      \item la velocidad del rotor a plena carga,
      \item la potencia en el eje del motor,
      \item el rendimiento.
    \end{enumerate}
  \item Un motor de inducción de 6 polos, 60 Hz, con una potencia
    absorbida de 100 kW, gira a 1150 R.P.M. Las pérdidas en el
    estator son de 300 W y las pérdidas mecánicas son de 350 W. Calcular:
    \begin{enumerate}
      \item el deslizamiento,
      \item las pérdidas en el rotor,
      \item la potencia útil en el eje del motor,
      \item el rendimiento.
    \end{enumerate}
  \item Un motor asincrónico trifásico de 10 polos, 2 CV, 400 V,
    conectado en estrella, gira a 68 R.P.M. El rendimiento a plena
    carga es del 80\% y el factor de potencia es de 0.8. Calcular:
    \begin{enumerate}
      \item la velocidad sincrónica,
      \item el deslizamiento a plena carga,
      \item la corriente de línea,
      \item la potencia entregada en el eje del motor.
    \end{enumerate}
  \item Un motor de inducción de 4 polos, 380 V, 50 Hz tiene una
    impedancia del estator de $z = 0.1 + j0.4 \Omega$. La impedancia
    del rotor en reposo es de $z = 0.12 + j0.3 \Omega$. El
    deslizamiento a plena carga es del 5\%. Calcular:
    \begin{enumerate}
      \item la corriente en el arranque,
      \item la corriente a plena carga,
      \item el par nominal si las pérdidas mecánicas son despreciables,
      \item el rendimiento si las pérdidas en el hierro son de 1000 W.
    \end{enumerate}
  \item Un motor asincrónico trifásico de 6 polos, conectado en
    estrella, con una potencia de 150 kW, gira a 1200 R.P.M. Las
    pérdidas en el rotor son de 500 W y las pérdidas en el estator
    son de 400 W. Calcular:
    \begin{enumerate}
      \item el deslizamiento,
      \item la potencia útil en el eje del motor,
      \item el rendimiento,
      \item el par desarrollado.
    \end{enumerate}
  \item Un motor de inducción trifásico de 4 polos, 440 V, 60 Hz,
    conectado en triángulo, tiene una impedancia del estator de $z =
    0.08 + j0.5 \Omega$. La impedancia del rotor en reposo es de $z =
    0.09 + j0.35 \Omega$. Calcular:
    \begin{enumerate}
      \item la corriente absorbida en el arranque,
      \item la corriente a plena carga con un deslizamiento del 2\%,
      \item la potencia nominal si las pérdidas mecánicas son despreciables,
      \item el rendimiento si las pérdidas en el hierro son de 1200 W.
    \end{enumerate}
  \item Un motor asincrónico trifásico de 8 polos, 2 CV, 380 V, tiene
    un deslizamiento de 1.5\% a plena carga. Las pérdidas en el
    estator y el rotor suman 600 W. Calcular:
    \begin{enumerate}
      \item la velocidad sincrónica,
      \item la velocidad del rotor a plena carga,
      \item la potencia entregada en el eje del motor,
      \item el rendimiento.
    \end{enumerate}
\end{enumerate}
\end{document}
